\documentclass[11pt]{article}

\usepackage[utf8]{inputenc}
\usepackage[T1]{fontenc}
\usepackage{lmodern}
\usepackage{geometry}
\usepackage{amsmath,amssymb,amsfonts,amsthm,mathtools}
\usepackage{bm}
\usepackage{enumitem}
\usepackage{microtype}
\usepackage{hyperref}
\usepackage{titlesec}
\usepackage{booktabs}
\usepackage{array}
\usepackage{setspace}
\usepackage{mathrsfs}
\usepackage{physics}

\geometry{margin=1in}
\hypersetup{colorlinks=true, linkcolor=black, urlcolor=blue, citecolor=black}
\setlist[enumerate]{topsep=0.4em,itemsep=0.35em}
\setlist[itemize]{topsep=0.3em,itemsep=0.25em}
\titleformat{\section}{\large\bfseries}{\thesection.}{0.5em}{}
\titleformat{\subsection}{\normalsize\bfseries}{\thesubsection.}{0.5em}{}
\setstretch{1.06}

\newcommand{\Aether}{\AE{}ther}
\newcommand{\Aflow}{\AE{}ther-flow}
\newcommand{\TheoryName}{\Aflow{} Interpretation of Relativity}
\newcommand{\TheoryProgram}{\Aether{} / \Aflow{} framework}
\newcommand{\ObserverMetric}{g^{\mathrm{br}}}
\newcommand{\CandidateMetric}{g^{\mathrm{cand}}}
\newcommand{\BaseShift}{\Sigma^{\mathrm{IER}}}
\newcommand{\ResponseShift}{\Sigma^{\mathrm{resp}}}
\newcommand{\SourceDensityRed}{\varrho_\Omega^{\mathrm{red}}}
\newcommand{\ReOff}{\widetilde{\mathcal R}_{\mu\nu}^{\mathrm{off}}}
\newcommand{\ReLongThree}{\widetilde{\mathcal R}_{\mu\nu}^{(\chi,3)}}
\newcommand{\ReLongFour}{\widetilde{\mathcal R}_{\mu\nu}^{(\chi,\ge 4)}}
\newcommand{\DeltaTCand}{\Delta T_{\mu\nu}^{\mathrm{cand}}}
\newcommand{\RemCand}{\mathcal R_{\mu\nu}^{\mathrm{cand}}}
\newcommand{\EinLin}{\mathcal E_{\ObserverMetric}}
\newcommand{\EinQuad}{\mathcal Q_{\ge 2}^{\mathrm{Ein}}}
\newcommand{\MixQuad}{\mathcal Q_{\mu\nu}^{\mathrm{mix}}}
\newcommand{\RespQuad}{\mathcal Q_{\mu\nu}^{\mathrm{resp}}}
\newcommand{\ResidualCore}{\mathcal V_{\mu\nu}^{\mathrm{res}}}
\newcommand{\ResidualDec}{\mathcal D_{\mu\nu}^{\mathrm{dec}}}
\newcommand{\GaugeExact}{\mathcal G_{\mu\nu}^{\mathrm{ex}}}
\newcommand{\ConstraintExact}{\mathcal C_{\mu\nu}^{\mathrm{ex}}}
\newcommand{\epsIR}{\varepsilon}

\newtheorem{definition}{Definition}
\newtheorem{proposition}{Proposition}
\newtheorem{remark}{Remark}
\newtheorem{corollary}{Corollary}
\newtheorem{theorem}{Theorem}

\title{\textbf{The \TheoryName{}}\\[0.4em]
\large Substrate-Response / Self-Response Charge-Polarization Candidate\\
\large Residual-Sector Verdict Note}
\author{}
\date{}

\begin{document}

\maketitle

\begin{abstract}
The benchmark-facing charge-polarization sequence has already fixed the concrete candidate, its full Phase D / E benchmark-gate verdicts, and its Phase F controlled GR-limit theorem. The immediate primary task is therefore no longer another packaging pass and no longer another anchored positive-pair continuation by default. It is Phase G: classify the finite-\(\epsIR\) remainder of that same candidate sector-by-sector and decide whether the branch is already fully closed at current scope.

The present note gives that verdict in the exact order now required. First, it rewrites the eliminated candidate remainder at fixed nonzero \(\epsIR\) into an exact finite-\(\epsIR\) decomposition
\[
\RemCand
=
\ResidualCore+\ResidualDec,
\]
where
\[
\ResidualCore
=
\ReLongFour
+\EinQuad[\BaseShift]
+\EinLin(\ResponseShift)
\]
is the only still-live quartic observer tensor on matter-free reduced data, while \(\ResidualDec\) consists of the mixed and pure response-quadratic pieces together with the one-metric matter re-expansion term. Second, it classifies every other leftover sector: the old off-longitudinal and cubic longitudinal blocks are absent on the candidate because they are already canceled by \(\BaseShift\); the inherited pure-gauge bookkeeping terms remain gauge exact; the eliminated source/polarization equations remain constraint exact; the mixed response-quadratic pieces and the matter re-expansion term are decoupled because they introduce no observer-accessible low-energy pole and vanish on the controlled infrared family.

The decisive point is that this verdict does \emph{not} rest only on the special exterior \(r^{-4}\) Hamiltonian neutralization. That earlier cancellation controls one projection of \(\ResidualCore\) on one static exterior regime, but it does not classify the full finite-\(\epsIR\) tensor away from that regime. At current disciplined scope, no theorem shows \(\ResidualCore\) to be absent, gauge exact, constraint exact, or decoupled on the full candidate branch. The first observer-accessible consequence is therefore only a quartic-order correction tensor in the finite-\(\epsIR\) observer equation, not a stronger phenomenological claim. The exact Phase G verdict is accordingly negative: the charge-polarization candidate is not yet fully closed at current scope and still leaves one genuine residual sector.
\end{abstract}

\section{Introduction}

The exact-closure benchmark remains the public front door of the repository, and the derivation-gates note remains the standard by which every bridge continuation is judged. The charge-polarization candidate note, the full infrared-spectrum / universal-coupling follow-up, and the controlled GR-limit note have already done the work they were supposed to do. The live question is now narrower than before.

That question is not whether the candidate has a controlled Einsteinian limit. It does. It is not whether the special exterior \(r^{-4}\) Hamiltonian density can be neutralized. It can. The Phase G question is sharper: at fixed nonzero \(\epsIR\), after the candidate has been eliminated honestly, which remainder sectors are actually absent, which are only gauge exact or constraint exact, which are merely heavy or one-metric re-expansion sectors that decouple, and whether anything genuinely observer-relevant survives after those classifications.

\paragraph{Framework and claim boundary.}
\TheoryProgram{} denotes the broader \Aether{} / \Aflow{} framework built on the claims that the \Aether{} is the underlying four-dimensional substrate of reality, the \Aflow{} is its intrinsic ordered motion, observed three-dimensional space is the local experiential slice of that deeper substrate, S-time is the experienced order of change arising from matter, light, and the \Aflow{}, and observed expansion is the three-dimensional appearance of deeper four-dimensional ordered motion. Within that framework, \TheoryName{} denotes the exact-closure theory in which the effective gravitational dynamics are adopted to be exactly Einsteinian with universal matter coupling. In the active sequence, \emph{adoption} means use of that established relativistic dynamics without claiming substrate derivation, while \emph{derivation} is reserved for a first-principles recovery from explicit substrate variables.

The present note stays strictly inside that boundary. It does not claim a completed finite-amplitude derivation of GR from substrate microphysics. It does not promote the current candidate to a fully closed deeper theory. Its narrower task is exactly the one fixed in the research plan and claim boundary: provide the residual-sector verdict for the recorded charge-polarization candidate.

\section{Fixed Phase G Inputs and Classification Classes}

\begin{definition}[Exact eliminated candidate remainder]
\label{def:remainder}
The recorded charge-polarization candidate is taken as a fixed input through its eliminated observer equation
\begin{equation}
G_{\mu\nu}[\CandidateMetric]
=
8\pi G_N\,T_{\mu\nu}^{\mathrm{matter}}[\CandidateMetric,\psi]
+\RemCand,
\label{eq:observer}
\end{equation}
with
\begin{equation}
\RemCand
=
\ReLongFour
+\EinLin(\ResponseShift)
+\EinQuad[\BaseShift+\ResponseShift]
-8\pi G_N\,\DeltaTCand.
\label{eq:remainder}
\end{equation}
\end{definition}

\begin{definition}[Phase G residual classes]
\label{def:classes}
A finite-\(\epsIR\) remainder block in \eqref{eq:remainder} is classified as follows.
\begin{enumerate}
    \item \textbf{Absent.} The block vanishes identically on the disciplined charge-polarization candidate branch and does not appear in the exact eliminated observer equation.
    \item \textbf{Gauge exact.} The block is a pure diffeomorphism or gauge identity after the same reduced gauge extraction already fixed by the benchmark-facing observer-remainder note.
    \item \textbf{Constraint exact.} The block is proportional to the eliminated source, polarization, or reduced slice equations and therefore vanishes on disciplined eliminated solutions.
    \item \textbf{Decoupled.} The block may be nonzero at fixed \(\epsIR\), but it is generated only by the already-screened elliptic/algebraic response variables or by one-metric matter re-expansion, it introduces no new observer-accessible infrared pole, and it tends to zero uniformly on every controlled infrared family of the Phase F note.
    \item \textbf{Genuine deviation sector.} A block is a genuine deviation sector at current scope if, after the previous four classes are removed, it remains as the leading finite-\(\epsIR\) observer-side correction and has not been reduced to any of those safer classes on the disciplined branch.
\end{enumerate}
\end{definition}

\begin{remark}
Definition \ref{def:classes} is operational rather than rhetorical. In particular, calling a block a genuine deviation sector here does not claim a full phenomenological prediction beyond current scope. It means only that the recorded manuscripts have not yet reduced that block to absence, exactness, or decoupling on the candidate branch.
\end{remark}

\begin{definition}[Fixed lower-order and linear-response identities]
\label{def:fixed}
The following earlier facts are treated as fixed inputs.
\begin{enumerate}
    \item The lower-order inverse-response shift obeys
    \begin{equation}
    \EinLin(\BaseShift)
    =
    -\ReOff
    -\ReLongThree,
    \label{eq:basesolve}
    \end{equation}
    so the old off-longitudinal and cubic longitudinal blocks are already canceled on the candidate branch.
    \item The added source/polarization sector is eliminated algebraically and elliptically; in particular, the polarization vector and trace-free tensor vanish, the scalar slots are fixed algebraically by the response potential, and
    \begin{equation}
    \delta\ResponseShift_{\mu\nu}=0
    \label{eq:linearzero}
    \end{equation}
    on the admissible homogeneous benchmark branch.
    \item On every controlled infrared family of the Phase F note,
    \begin{equation}
    \BaseShift=\mathcal O(\epsIR^2),
    \qquad
    \ResponseShift=\mathcal O(\epsIR^4),
    \qquad
    \ReLongFour=\mathcal O(\epsIR^4),
    \qquad
    \DeltaTCand=\mathcal O(\epsIR^4).
    \label{eq:phaseFbounds}
    \end{equation}
\end{enumerate}
\end{definition}

\section{Exact Finite-\texorpdfstring{\(\epsIR\)}{epsilon} Decomposition}

\begin{definition}[Mixed and pure response-quadratic sectors]
\label{def:split}
Define the exact bilinear and pure response-quadratic pieces of \(\EinQuad[\BaseShift+\ResponseShift]\) by
\begin{equation}
\EinQuad[\BaseShift+\ResponseShift]
=
\EinQuad[\BaseShift]
+\MixQuad
+\RespQuad,
\label{eq:quadsplit}
\end{equation}
where \(\MixQuad\) is bilinear in \(\BaseShift\) and \(\ResponseShift\), and \(\RespQuad\) is at least quadratic in \(\ResponseShift\).
\end{definition}

\begin{proposition}[Exact remainder decomposition away from the exterior regime]
\label{prop:exactsplit}
For the eliminated charge-polarization candidate,
\begin{equation}
\RemCand
=
\ResidualCore+\ResidualDec,
\label{eq:exactsplit}
\end{equation}
where
\begin{equation}
\ResidualCore
\equiv
\ReLongFour
+\EinQuad[\BaseShift]
+\EinLin(\ResponseShift),
\label{eq:core}
\end{equation}
and
\begin{equation}
\ResidualDec
\equiv
\MixQuad+\RespQuad-8\pi G_N\,\DeltaTCand.
\label{eq:dec}
\end{equation}
This decomposition is exact at fixed nonzero \(\epsIR\) and does not use the decisive exterior cancellation regime.
\end{proposition}

\begin{proof}
Substitute the split \eqref{eq:quadsplit} into the exact candidate remainder formula \eqref{eq:remainder}. The result is \eqref{eq:exactsplit}--\eqref{eq:dec} with no appeal to asymptotic exterior arguments, no use of \(r^{-4}\) neutralization, and no use of the \(\epsIR\to0\) limit.
\end{proof}

\begin{remark}
Equation \eqref{eq:exactsplit} is the right Phase G starting point. The residual-sector verdict must classify the exact finite-\(\epsIR\) pieces of \(\RemCand\), not only the limiting statement that \(\RemCand=\mathcal O(\epsIR^4)\).
\end{remark}

\section{Sector-by-Sector Classification}

\begin{proposition}[Absent, gauge-exact, and constraint-exact sectors]
\label{prop:absent}
At current candidate scope, the following sector classifications hold.
\begin{enumerate}
    \item \textbf{Absent lower-order sectors.} The old off-longitudinal block \(\ReOff\) and the first cubic longitudinal block \(\ReLongThree\) are absent from the exact eliminated candidate remainder.
    \item \textbf{Gauge-exact bookkeeping sector.} Any inherited pure-gauge bookkeeping terms remain gauge exact and do not survive as an independent block in \eqref{eq:remainder}.
    \item \textbf{Constraint-exact elimination sector.} Terms proportional to the eliminated source/polarization equations and to the reduced slice equations remain constraint exact and vanish on disciplined eliminated solutions.
    \item \textbf{Absent linear response sector.} No unsuppressed linear response block survives: the added response sector is absent at linear observer-equation level.
\end{enumerate}
\end{proposition}

\begin{proof}
Item (1) is exactly the fixed lower-order identity \eqref{eq:basesolve}: once \(\BaseShift\) is kept fixed, \(\ReOff\) and \(\ReLongThree\) no longer appear in \eqref{eq:remainder}. Item (2) is inherited from the same reduced gauge extraction used by the earlier observer-remainder classification note; the present candidate note writes the remainder only after that gauge bookkeeping has already been removed. Item (3) follows from the Euler--Lagrange elimination of the charge, multiplier, and polarization variables together with the already-recorded reduced slice package: any term still proportional to those equations vanishes on the eliminated branch. Item (4) is \eqref{eq:linearzero}, which shows that the charge-polarization sector does not generate a new linear observer imprint.
\end{proof}

\begin{proposition}[Decoupled sectors]
\label{prop:decoupled}
The sectors \(\MixQuad\), \(\RespQuad\), and \(\DeltaTCand\) are decoupled in the sense of Definition \ref{def:classes}.
\end{proposition}

\begin{proof}
By Definition \ref{def:split}, \(\MixQuad\) is bilinear in \(\BaseShift\) and \(\ResponseShift\), while \(\RespQuad\) is at least quadratic in \(\ResponseShift\). Using \eqref{eq:phaseFbounds},
\begin{equation}
\MixQuad=\mathcal O(\epsIR^6),
\qquad
\RespQuad=\mathcal O(\epsIR^8)
\label{eq:highorder}
\end{equation}
on every controlled infrared family. These two blocks are therefore generated only by the already-screened heavy response sector and are at least two reduced-data orders below the leading quartic observer correction.

For \(\DeltaTCand\), the full Phase D / E gate note already fixed that all matter and light sectors couple only through the single operative metric \(\CandidateMetric\). The term \(\DeltaTCand\) is therefore not an independent second-source law but only the re-expansion of ordinary matter stress under the same one operative metric. It vanishes on matter-free reduced data and satisfies \(\DeltaTCand=\mathcal O(\epsIR^4)\) on every controlled infrared family by \eqref{eq:phaseFbounds}. Since the same Phase D note already proved that the candidate introduces no new observer-accessible linear pole, scalar, vector, preferred-frame, ghost, tachyonic, or second-cone sector, \(\DeltaTCand\) is decoupled rather than a new observer-accessible branch.
\end{proof}

\begin{remark}
Proposition \ref{prop:decoupled} is the precise sense in which the Phase G verdict goes beyond the controlled limit without collapsing back into it. These sectors are classified at fixed nonzero \(\epsIR\) by their structural origin, while the Phase F estimates certify that they really do decouple on the controlled family.
\end{remark}

\section{Why the Verdict Cannot Rest Only on the Exterior \texorpdfstring{\(r^{-4}\)}{r^-4} Neutralization}

\begin{proposition}[Special exterior neutralization controls only one projection]
\label{prop:outside}
On the decisive rotational matter-free exterior regime, the earlier candidate note proves only that
\begin{equation}
n^\mu n^\nu \ResidualCore
=
\mathcal O(r^{-5}).
\label{eq:extneutral}
\end{equation}
This does not classify the full finite-\(\epsIR\) tensor \(\ResidualCore\) away from that special regime.
\end{proposition}

\begin{proof}
The special exterior theorem of the candidate note applies on a static exterior region with the Coulomb tail and its Schwarzschild-type response dressing isolated explicitly. On matter-free reduced data, the quartic remainder there is exactly \(\ResidualCore\), because the mixed and pure response-quadratic pieces are higher order and \(\DeltaTCand\) vanishes. The recorded neutralization cancels the leading \(r^{-4}\) Hamiltonian projection of that quartic tensor, yielding \eqref{eq:extneutral}.

But \eqref{eq:extneutral} is only a statement about the contraction \(n^\mu n^\nu \ResidualCore\) on one asymptotic exterior regime. It does not determine the full tensor components of \(\ResidualCore\), it does not classify interior or nonstatic regions, and it does not convert \(\ResidualCore\) into a gauge identity, a constraint identity, or a decoupled heavy block. Therefore the Phase G verdict cannot rest only on \eqref{eq:extneutral}.
\end{proof}

\begin{corollary}[Only one still-live sector remains]
\label{cor:oneleft}
After Propositions \ref{prop:absent} and \ref{prop:decoupled}, the only still-live finite-\(\epsIR\) candidate block is \(\ResidualCore\).
\end{corollary}

\begin{proof}
Proposition \ref{prop:absent} removes the absent, gauge-exact, and constraint-exact sectors. Proposition \ref{prop:decoupled} removes \(\ResidualDec\). By the exact decomposition \eqref{eq:exactsplit}, only \(\ResidualCore\) remains.
\end{proof}

\section{First Observer-Accessible Consequence and Phase G Verdict}

\begin{proposition}[First observer-accessible consequence if a residual sector survives]
\label{prop:consequence}
Suppose \(\ResidualCore\) is not reduced to absence, gauge exactness, constraint exactness, or decoupling on the disciplined candidate branch. Then the first observer-accessible consequence at current scope is a quartic-order correction tensor in the finite-\(\epsIR\) observer equation:
\begin{equation}
G_{\mu\nu}[\CandidateMetric_{\epsIR}]
=
8\pi G_N\,T_{\mu\nu}^{\mathrm{matter}}[\CandidateMetric_{\epsIR},\psi_{\epsIR}]
+\epsIR^4\,\mathfrak V_{\mu\nu}^{(\epsIR)}
+\epsIR^4\,\mathfrak D_{\mu\nu}^{(\epsIR)},
\label{eq:firstcons}
\end{equation}
where
\begin{equation}
\mathfrak V_{\mu\nu}^{(\epsIR)}
\equiv
\epsIR^{-4}{\ResidualCore}^{(\epsIR)}
\label{eq:vdef}
\end{equation}
and
\begin{equation}
\mathfrak D_{\mu\nu}^{(\epsIR)}
\equiv
\epsIR^{-4}\ResidualDec
\label{eq:ddef}
\end{equation}
on any controlled infrared family for which \(\epsIR>0\). No stronger phenomenological claim is justified at current scope.
\end{proposition}

\begin{proof}
Equation \eqref{eq:firstcons} is just the exact observer equation \eqref{eq:observer} with the exact decomposition \eqref{eq:exactsplit} and the rescaled notation \eqref{eq:vdef}--\eqref{eq:ddef}. The Phase F bounds show that \(\ResidualCore=\mathcal O(\epsIR^4)\) and \(\ResidualDec=\mathcal O(\epsIR^4)\) on the controlled family, while Proposition \ref{prop:decoupled} shows that \(\ResidualDec\) is structurally decoupled. Proposition \ref{prop:outside} shows that the existing manuscripts do not identify \(\ResidualCore\) with a safer class away from the special exterior projection. Hence the first honest observer-side consequence is only the quartic correction tensor \(\epsIR^4\mathfrak V_{\mu\nu}^{(\epsIR)}\) itself.
\end{proof}

\begin{table}[t]
\centering
\begin{tabular}{p{0.35\textwidth}p{0.18\textwidth}p{0.36\textwidth}}
\toprule
\textbf{Sector} & \textbf{Verdict} & \textbf{Reason at current scope} \\
\midrule
Old \(\ReOff\) and \(\ReLongThree\) blocks & absent & Already canceled by the fixed lower-order inverse-response shift \(\BaseShift\). \\
Inherited pure-gauge bookkeeping terms & gauge exact & Removed by the same reduced gauge extraction used before the candidate remainder is written. \\
Eliminated source / polarization / reduced-slice equations & constraint exact & Vanish on disciplined eliminated solutions after the charge-polarization Euler--Lagrange reduction. \\
Standalone linear response block & absent & \(\delta\ResponseShift_{\mu\nu}=0\), so no unsuppressed linear observer imprint survives. \\
\(\MixQuad\) and \(\RespQuad\) & decoupled & Generated only by the already-screened elliptic/algebraic response block; \(\mathcal O(\epsIR^6)\) and \(\mathcal O(\epsIR^8)\) on the controlled family. \\
\(\DeltaTCand\) & decoupled & One-metric matter re-expansion term, not a second matter source; vanishes on matter-free data and no new low-energy pole is introduced. \\
\(\ResidualCore=\ReLongFour+\EinQuad[\BaseShift]+\EinLin(\ResponseShift)\) & genuine deviation sector & Only one exterior Hamiltonian projection is neutralized; away from that regime the full finite-\(\epsIR\) tensor is not yet absent, exact, or decoupled. \\
\bottomrule
\end{tabular}
\caption{Phase G sector classification for the charge-polarization candidate.}
\label{tab:phaseg}
\end{table}

\begin{theorem}[Phase G residual-sector verdict]
\label{thm:phaseg}
At the present disciplined scope, the charge-polarization candidate is \emph{not} fully closed. More precisely:
\begin{enumerate}
    \item every leftover finite-\(\epsIR\) block other than \(\ResidualCore\) is classified by Table \ref{tab:phaseg} as absent, gauge exact, constraint exact, or decoupled;
    \item the special exterior \(r^{-4}\) Hamiltonian neutralization does not by itself classify the full tensor \(\ResidualCore\) away from that regime;
    \item therefore \(\ResidualCore\) is the only still-live genuine residual sector at current scope;
    \item the exact candidate verdict is negative: the recorded charge-polarization branch still leaves a genuine residual sector and cannot yet be promoted as a fully closed derivation of GR.
\end{enumerate}
\end{theorem}

\begin{proof}
Item (1) is exactly Propositions \ref{prop:absent} and \ref{prop:decoupled}, summarized in Table \ref{tab:phaseg}. Item (2) is Proposition \ref{prop:outside}. Item (3) is Corollary \ref{cor:oneleft}: once the absent, exact, and decoupled sectors are removed, only \(\ResidualCore\) remains. By Definition \ref{def:classes}, that surviving leading finite-\(\epsIR\) observer-side correction is therefore a genuine deviation sector at current scope. Item (4) is then the claim-boundary verdict: full closure would require \(\ResidualCore\) to be absent, gauge exact, constraint exact, or decoupled as well, and no such result is presently recorded.
\end{proof}

\begin{corollary}[What should be next]
The next primary manuscript should not be another packaging pass and should not default to another anchored positive-pair continuation. The next primary work should instead be a deeper observer-equation redesign or equivalent substrate-response completion that removes or reclassifies \(\ResidualCore\), unless some side line actually changes that observer equation.
\end{corollary}

\section{Conclusion}

The active next task after the controlled GR-limit note was Phase G itself: the residual-sector verdict for the same charge-polarization candidate. The present manuscript performs that step at the exact scope now required. It classifies the finite-\(\epsIR\) remainder sector-by-sector, not only in the \(\epsIR\to0\) limit, and it does so away from the decisive exterior cancellation regime rather than resting only on the special \(r^{-4}\) Hamiltonian neutralization.

The positive result is a real narrowing of the burden. The old off-longitudinal and cubic longitudinal sectors are absent on the candidate, the inherited gauge bookkeeping remains gauge exact, the eliminated source/polarization equations remain constraint exact, and the mixed response-quadratic plus one-metric matter re-expansion sectors are decoupled. The first observer-accessible consequence is therefore not a new linear pole, not a second matter cone, and not a preferred-frame matter law. If a deviation survives, it survives first as a quartic-order correction tensor in the finite-\(\epsIR\) observer equation.

The exact verdict is therefore sharp. At current disciplined scope, the charge-polarization candidate is not fully closed. The only still-live sector is the composite quartic tensor \(\ResidualCore=\ReLongFour+\EinQuad[\BaseShift]+\EinLin(\ResponseShift)\), and away from the special exterior neutralization regime it has not yet been reduced to absence, gauge exactness, constraint exactness, or decoupling. The candidate still leaves a genuine residual sector. The next honest primary work is to remove or reclassify that observer-equation sector itself, not another packaging pass and not another anchored positive-pair continuation unless that side line changes the observer equation.

\end{document}
